\section{Tekis o‘zgaruvchan harakat. Tezlanish}

\begin{enumerate}

\item Jism tekis tezlanuvchan harakatlanib, 6 s davomida o‘z tezligini 0,1 m/s dan 0,4 m/s gacha oshirdi. 
Jism necha m/s$^2$ tezlanish bilan harakatlangan?

\begin{enumerate}[label=\Alph*)]
\item 5
\item 0,5
\item 0,05
\item 0,005
\end{enumerate}

\item To‘g‘ri chiziqli tekis sekinlanuvchan harakatda tezlanish vektori ... 

\begin{enumerate}[label=\Alph*)]
\item moduli o‘zgarmas bo‘lib, yo‘nalishi uzluksiz o‘zgarib turadi.
\item moduli va yo‘nalishi o‘zgarib turadi.
\item  $a>0$ holda moduli va yo‘nalishi o‘zgarmaydi. 
\item  $a<0$ holda moduli va yo‘nalishi o‘zgarmaydi.
\end{enumerate}

\item Quyidagi gapni to‘g‘ri mazmunda to‘ldiruvchi javobni belgilang:

Tekis o‘zgaruvchan harakatdagi oniy tezlikni hisoblash uchun ... berilishi lozim.

\begin{enumerate}[label=\Alph*)]
\item harakatning boshlang‘ich tezligi va tezlanishi
\item tezlanish va harakatlanish vaqti
\item boshlang‘ich tezlik va tezlanishning moduli
\item tezlanish, harakatlanish vaqti va boshlang‘ich tezlik
\end{enumerate}

\item Agar avtomobilning boshlang‘ich tezligi 36 km/h bo‘lib, 1 minutdan so‘ng to‘xtasa, tormozlanish jarayonida tezlikning vaqtga bog‘lanish tenglamasi qanday bo‘ladi?

\begin{enumerate}[label=\Alph*)]
\item $v=10-20t$
\item $v=10-0,17t$
\item $v=36-36t$
\item $v=36+36t$
\end{enumerate}

\item Jism $v_0$ boshlang‘ich tezlik va $a$ tezlanish bilan tekis
sekinlanuvchan harakat qilmoqda. Jismning tezligi qancha vaqtdan
so‘ng 2 marta kamayadi?

\begin{enumerate}[label=\Alph*)]
\item $v_0/a$
\item $2v_0/a$
\item $v_0/3a$
\item $v_0/2a$
\end{enumerate}

\item Jism tinch holatdan harakatga kelib, $t$ vaqtda $2a$ tezlanish
bilan, so‘ngra $4t$ vaqtda tekis harakat qildi. Oxirida yana $2t$
vaqtda $a$ tezlanish bilan harakat qildi. Jismning oxirgi tezligini toping.

\begin{enumerate}[label=\Alph*)]
\item $-at$
\item $at$
\item $4at$
\item $0$
\end{enumerate}



\item 20 m/s tezlik bilan ketayotgan avtomobil tormozlanish
natijasida 2 s da to‘xtadi. Avtomobilning tormozlanish
yo‘lini aniqlang (m).

\begin{enumerate}[label=\Alph*)]
\item 200
\item 150
\item 100
\item 50
\end{enumerate}

\item Boshlang‘ich tezlik bilan to‘g‘ri chiziqli tekis
tezlanuvchan harakat qilayotgan jismning boshlang‘ich
tezligi va tezlanishi 3 marta oshirilsa, uning birlik vaqt
ichidagi ko‘chishi necha marta ortadi?

\begin{enumerate}[label=\Alph*)]
\item 1,5
\item 4,5
\item 3
\item 6
\end{enumerate}

\item Quyidagi tenglamalarning qaysilari tekis harakatga taalluqli:

1) $s=2t$ \quad
2) $s=at^2$ \quad
3) $x=3t+2$ \quad
4) $v=4+t$ \quad
5) $v=5$

\begin{enumerate}[label=\Alph*)]
\item 2,3,4
\item 1,2,3
\item 1,3,5
\item 3,4,5
\end{enumerate}

\item Moddiy nuqta tezligining vaqt bo‘yicha o‘zgarish qonuni
$v=4t$ ko‘rinishga ega. Uning 5 sekundda o‘tgan yo‘li qanday (m)?

\begin{enumerate}[label=\Alph*)]
\item 32
\item 16
\item 18
\item 50
\end{enumerate}

\item Boshlang‘ich tezliksiz tekis tezlanuvchan harakatlanayotgan jism 
harakatning 5-sekundida 27 m yo‘l bosgan bo‘lsa, uning tezlanishi qanday (m/s$^2$)?

\begin{enumerate}[label=\Alph*)]
\item 3,5
\item 5,4
\item 6
\item 9
\end{enumerate}


\item Reaktiv samolyotga yerda 640 m uzunlikdagi uchish yo‘li talab qilinadi. 
Uning motorlari unga 5 m/s$^2$ tezlanish bera oladi. Samolyot faqat o‘z 
motorlari yordamida tezligi 72 km/h bo‘lgan aviatashuvchi kemadan 
uchishi uchun kema palubasining uzunligi qanday bo‘lishi kerak?

\begin{enumerate}[label=\Alph*)]
\item 64
\item 360
\item 400
\item 480
\end{enumerate}


\item Jism harakatni boshlab, harakat vaqtining boshlang‘ich 
$1/3$ qismida tekis tezlanuvchan, so‘ngra tekis va oxirgi 
$1/4$ qismida tekis sekinlanuvchan harakatlanib to‘xtadi. 

Bu harakatda jismning o‘rtacha tezligi $17\ \text{m/s}$ ga teng 
bo‘lgan bo‘lsa, jism tezligining eng katta qiymati qanday (m/s)?

\begin{enumerate}[label=\Alph*)]
\item 24
\item 20
\item 34
\item 51
\end{enumerate}


\item Jismning harakat tenglamasi $x = 3 + 8t - t^2$ ko‘rinishga ega. Jism tezligi nolga teng bo‘lguncha qanday (m) yo‘l o‘tadi?

\begin{enumerate}[label=\Alph*)]
\item 3
\item 8
\item 16
\item 19
\end{enumerate}


\item Moddiy nuqtaning harakat tenglamasi $x = 12t - 2t^2 $ko‘rinishga ega. U $t=0$ paytdan boshlab, to‘xtaguncha 
qanday (m) yo‘l bosadi?

\begin{enumerate}[label=\Alph*)]
\item 12
\item 18
\item 16
\item 32
\end{enumerate}

\end{enumerate}



\begin{table}[h!]
\centering
\begin{tabular}{|c|c|c|c|c|c|c|c|c|c|}
\hline
1c & 2d & 3d & 4b & 5d & 6c & 7d & 8c & 9c & 10d \\ \hline
11b & 12d & 13c & 14c & 15b &  &  &  &  &  \\ \hline
\end{tabular}
\end{table}










